\documentclass[12pt,a4paper]{article}
\usepackage[utf8]{inputenc}
\usepackage[vietnamese]{babel}
\usepackage[left=2cm,right=2cm,top=2cm,bottom=2cm]{geometry}
\usepackage{mathptmx}
\usepackage{enumitem}

\pagestyle{plain}
\linespread{1.15}

\begin{document}

% --- HEADER QUỐC HIỆU ---
\noindent
\begin{minipage}[t]{0.45\textwidth}
    \begin{center}
        \small
        CÔNG AN THÀNH PHỐ HÀ NỘI\\
        \textbf{\underline{PHÒNG CẢNH SÁT HÌNH SỰ}}\\[0.1cm]
        Số: 11/BB-TXCAM\\
        \textit{V/v trích xuất dữ liệu camera và kỷ vật trốn tìm}
    \end{center}
\end{minipage}
\hfill
\begin{minipage}[t]{0.52\textwidth}
    \begin{center}
        \small
        \textbf{CỘNG HÒA XÃ HỘI CHỦ NGHĨA VIỆT NAM}\\
        \textbf{\underline{Độc lập - Tự do - Hạnh phúc}}\\[0.1cm]
        \textit{Hà Nội, ngày 25 tháng 07 năm 2026}
    \end{center}
\end{minipage}

\vspace{0.6cm}

\begin{center}
    \textbf{\large BIÊN BẢN TRÍCH XUẤT DỮ LIỆU CAMERA VÀ VẬT CHỨNG}\\
    \textit{(Phân tích lịch trình di chuyển của nghi phạm Tùng)}
\end{center}

\vspace{0.3cm}

\section*{I. KẾT QUẢ TRÍCH XUẤT CAMERA CÂY XĂNG (KÝ HIỆU f3-1)}
\begin{itemize}[leftmargin=1.5cm]
    \item \textbf{Địa điểm camera:} Cây xăng Petrolimex số 12, cách ngõ nhà Khang 300m.
    \item \textbf{Lịch trình nghi phạm Tùng (đi xe máy Honda Wave đen, BKS 29H1-829.12):}
    \begin{itemize}[leftmargin=0.8cm]
        \item \textbf{19:45:} Tùng rẽ xe máy vào hẻm dẫn tới nhà Khang.
        \item \textbf{20:15:} Camera ghi nhận Tùng **chạy xộc ra hoảng loạn**, không đội mũ bảo hiểm, nhảy lên xe rồ ga tháo chạy tốc độ cao.
    \end{itemize}
\end{itemize}

\section*{II. TÀI LIỆU KỶ VẬT TRÒ TRỐN TÌM NĂM 1998 (KÝ HIỆU n3, p4, p5)}
\begin{itemize}[leftmargin=1.5cm]
    \item \textbf{Mảnh giấy note `n3`:} Ghi chép vị trí trốn của 4 đứa trẻ năm 1998. Dòng chữ viết tay của Khang: *"Gia Huy núp trong khoang tủ gỗ âm tường"*.
    \item \textbf{Ảnh chụp `p4`:} 4 đứa trẻ (Khang, Tùng, Mai, Gia Huy) chơi trốn tìm mùa hè năm 1998.
    \item \textbf{Bài báo cũ `p5`:} Báo Công an Thủ đô ra ngày 12/07/1998 đưa tin bé trai Nguyễn Gia Huy (6 tuổi) tử vong do ngạt khí trong tủ gỗ đóng chốt ngoài.
\end{itemize}

\vspace{0.8cm}

\noindent
\begin{minipage}[t]{0.45\textwidth}
    \small
    \textbf{\textit{Nơi nhận:}}\\
    - Trưởng Ban Chuyên án;\\
    - Lưu: VT, CSHS.
\end{minipage}
\hfill
\begin{minipage}[t]{0.5\textwidth}
    \begin{center}
        \small
        \textbf{CÁN BỘ TRÍCH XUẤT}\\
        \textit{(Ký, ghi rõ họ tên)}\\[1.5cm]
        \textbf{Trung úy Phạm Thanh}
    \end{center}
\end{minipage}

\end{document}
